% =============================================================================
%  Professional Thesis report - IMT Nord Europe (DNM)
%
%  Single-file framework, built for Overleaf: upload this file plus
%  references.bib and any images. Nothing else is \input.
%
%  Structure and limits follow the supplied instructions
%  "230614_FOR_DNM Professional Thesis Report Instructions.pdf", kept outside
%  this repository.
%
%  DRAFT STATUS. Technical content in the Introduction, Section 2, and the
%  Conclusion is drafted from the project's working diaries and dated plans
%  (03/04/2026 to 03/08/2026) and preserves their verified / prepared / unrun
%  distinctions.
%  Company-specific and personal content (Section 1, Section 3, cover fields,
%  acknowledgements) is left as TODO because it cannot be derived from the
%  repository and must not be invented.
%
%  PRESENTATION SETTINGS ARE PROVISIONAL. The instructions defer layout to a
%  separate document derived from experimental standard Z41 006, obtainable
%  through IMT My Services. Confirm fonts, margins, and spacing against it
%  before final submission. Erasmus+ does not prescribe a universal final-
%  report font or page layout, so no generic Erasmus template is applied.
%
%  Navigate this file by searching for the "== " banners.
% =============================================================================

\documentclass[a4paper,12pt,oneside,openany]{report}

\usepackage[T1]{fontenc}
\usepackage[utf8]{inputenc}
\usepackage[french,english]{babel}
\usepackage{lmodern}
\usepackage{microtype}
\usepackage[left=2.5cm,right=2.5cm,top=1.5cm,bottom=2cm]{geometry}

\usepackage{graphicx}
% Overleaf: keep images in a figures/ folder inside the project. Do not point
% this at a sibling directory - those paths do not exist once uploaded.
\graphicspath{{figures/}}

\usepackage{booktabs}
\usepackage{array}
\usepackage{tabularx}
\usepackage{longtable}

% Ragged-right column types. Narrow justified p{} columns force TeX to stretch
% interword space, which is what produced every "Underfull \hbox" warning on the
% first compile. Ragged-right removes the cause rather than tolerating it, and
% reads better in narrow table cells.
\newcolumntype{L}[1]{>{\raggedright\arraybackslash}p{#1}}
\newcolumntype{Y}{>{\raggedright\arraybackslash}X}
\usepackage{caption}
\usepackage{chngcntr}
\usepackage{xcolor}
\usepackage{tikz}
\usepackage[nottoc]{tocbibind}
\usepackage[hidelinks]{hyperref}

% The instructions require figures and tables to carry UNIQUE ASCENDING ARABIC
% numbers throughout, including in the appendices. The report class would
% otherwise scope both counters to the chapter, producing "Table 0.1" in the
% front matter and "Figure A.1" in the appendix - neither is Arabic-continuous.
\counterwithout{figure}{chapter}
\counterwithout{table}{chapter}

% NOTE: these two lines only tell the caption package how to SPACE a caption.
% They do not move it. Placement comes from where you write \caption:
%   figure -> \caption AFTER \includegraphics  (caption below, as required)
%   table  -> \caption BEFORE the tabular      (caption above, as required)
\captionsetup[figure]{position=bottom}
\captionsetup[table]{position=top}

\setcounter{tocdepth}{2}
\setcounter{secnumdepth}{3}
\setlength{\emergencystretch}{2em}

\newcommand{\ReportTodo}[1]{\textbf{TODO:} #1}

% == FIGURE PLACEHOLDERS ======================================================
% \PlaceholderFigure{label-suffix}{what the figure will show}{caption}
% Draws a captioned, numbered, referenceable grey box so the document compiles
% and the list of figures is correct before any image exists. Replace the box
% with \includegraphics{...} as each image is produced; keep caption and label.
% Caption sits AFTER the box, so it prints below the figure, as required.
\newcommand{\PlaceholderFigure}[3]{%
  \begin{figure}[tbp]
    \centering
    \fcolorbox{black!40}{black!5}{%
      \begin{minipage}[c][0.22\textheight][c]{0.88\textwidth}
        \centering
        {\itshape\color{black!60} [FIGURE PLACEHOLDER]\par\medskip #2}
      \end{minipage}}
    \caption{#3}
    \label{fig:#1}
  \end{figure}}

% == COVER FIELDS =============================================================
% Required by the instructions: student name, graduating class, company name,
% type of Professional Thesis, and a title or overview of the subject.
\newcommand{\StudentName}{[Student name]}
\newcommand{\GraduatingClass}{[Graduating class]}
\newcommand{\CompanyName}{[Host institution]}
\newcommand{\ThesisType}{DNM Professional Thesis}
\newcommand{\ReportTitle}{Autonomous Surface Vehicle for Water-Quality Survey:
  Simulation, Onboard Perception, and Field Bring-Up}
\newcommand{\SubjectOverview}{Development and evidence-based evaluation of a ROS 2 navigation,
  perception, and telemetry foundation for an autonomous surface vehicle}

% Optional fields: keep only those approved for the final cover.
\newcommand{\InstitutionName}{IMT Nord Europe}
\newcommand{\CompanySupervisor}{[Company supervisor]}
\newcommand{\AcademicSupervisor}{[Academic supervisor]}
\newcommand{\ReportDate}{\today}

% Set \reportconfidentialtrue ONLY after the host company or internship
% supervisor confirms it. A confidential report is NOT submitted to
% MyLearningSpace - see README.md for the separate delivery route.
\newif\ifreportconfidential
\reportconfidentialfalse

\hypersetup{
  pdftitle={\ReportTitle},
  pdfauthor={\StudentName},
  pdfsubject={Professional Thesis Report}
}

\begin{document}
\selectlanguage{english}

% == COVER PAGE ===============================================================
\begin{titlepage}
  \centering
  {\Large \InstitutionName\par}
  \vspace{1.5cm}
  {\large \ThesisType\par}
  \vspace{1cm}
  {\Huge\bfseries \ReportTitle\par}
  \vspace{0.8cm}
  {\Large \SubjectOverview\par}
  \vfill

  \begin{tabular}{@{}ll@{}}
    \textbf{Student:} & \StudentName \\
    \textbf{Graduating class:} & \GraduatingClass \\
    \textbf{Host institution:} & \CompanyName \\
    \textbf{Company supervisor (optional):} & \CompanySupervisor \\
    \textbf{Academic supervisor (optional):} & \AcademicSupervisor \\
    \textbf{Date:} & \ReportDate
  \end{tabular}

  \ifreportconfidential
    \vfill
    {\Large\bfseries CONFIDENTIAL\par}
  \fi
\end{titlepage}

\clearpage
\pagenumbering{roman}

% == ACKNOWLEDGEMENTS =========================================================
\chapter*{Acknowledgements}
\addcontentsline{toc}{chapter}{Acknowledgements}

\ReportTodo{Thank the people and organizations that materially supported the internship and
the Professional Thesis. Confirm names, roles, and preferred forms of address before
submission.}

% == CONTENTS AND LISTS =======================================================
\tableofcontents
\listoffigures
\listoftables

% == ENGLISH/FRENCH KEYWORD GLOSSARY ==========================================
\chapter*{English/French Keyword Glossary}
\addcontentsline{toc}{chapter}{English/French Keyword Glossary}

\ReportTodo{Review the French equivalents marked ``confirm'' with a French-speaking supervisor
before submission; the rest are standard.}

Table~\ref{tab:glossary-en-fr} lists the English and French keywords used throughout the report.

\begin{longtable}{@{}L{0.27\textwidth}L{0.30\textwidth}L{0.33\textwidth}@{}}
  \caption{English/French glossary of the report's key words}\label{tab:glossary-en-fr}\\
  \toprule
  \textbf{English keyword} & \textbf{French keyword} & \textbf{Note} \\
  \midrule
  \endfirsthead
  \caption[]{English/French glossary of the report's key words (continued)}\\
  \toprule
  \textbf{English keyword} & \textbf{French keyword} & \textbf{Note} \\
  \midrule
  \endhead
  Unmanned surface vehicle (USV) & véhicule de surface autonome & --- \\
  Digital twin & jumeau numérique & standard \\
  Autonomous navigation & navigation autonome & --- \\
  Obstacle avoidance & évitement d'obstacles & --- \\
  Path planning & planification de trajectoire & --- \\
  Waypoint & point de passage & --- \\
  Coverage path & trajectoire de couverture (boustrophédon) & --- \\
  Real-time factor & facteur temps réel & confirm; often left in English \\
  Flight controller & contrôleur de vol / pilote automatique & --- \\
  Telemetry & télémétrie & --- \\
  Companion computer & ordinateur embarqué compagnon & confirm \\
  Edge AI accelerator & accélérateur d'IA embarquée & confirm \\
  Object detection & détection d'objets & --- \\
  Water-quality monitoring & surveillance de la qualité de l'eau & --- \\
  Cellular automaton & automate cellulaire & standard \\
  \bottomrule
\end{longtable}

% == OVERALL PROJECT SCHEDULE =================================================
\chapter*{Overall Project Schedule}
\addcontentsline{toc}{chapter}{Overall Project Schedule}

The schedule below covers recorded work from 03/04/2026 to 01/08/2026 and the prepared plan through
03/08/2026. Phases are colour-coded by evidence state: \textbf{verified} work was executed and
observed, while \textbf{prepared} work was built and source-checked but never run. The internship
continued past the right-hand edge of the chart, so the schedule is a progress record rather than a
completed plan.

Figure~\ref{fig:gantt} summarizes the sequence and overlap of the recorded phases, while
Table~\ref{tab:project-schedule} gives their exact dates and evidence states.

\begin{figure}[tbp]
  \centering
  \begin{tikzpicture}[x=0.085cm, y=-0.62cm, font=\scriptsize]
    % ---- month grid: day 0 = 03/04/2026 -------------------------------------
    \foreach \dx/\m in {0/April, 28/May, 59/June, 89/July, 120/Aug.}{
      \draw[black!20] (\dx,0.3) -- (\dx,15.0);
      \node[anchor=south, black!65] at (\dx,0.3) {\m};
    }
    \draw[black!20] (122,0.3) -- (122,15.0);

    % ---- bars: {row}{start-day}{end-day}{colour}{label} ---------------------
    \def\gbar#1#2#3#4#5{%
      \fill[#4, rounded corners=1pt] (#2,{#1-0.26}) rectangle (#3,{#1+0.26});
      \node[anchor=east, black!85] at (-3,#1) {#5};}

    \definecolor{cver}{RGB}{31,110,165}
    \definecolor{cpre}{RGB}{214,158,46}

    \gbar{1}{0}{6}{cver}{Code review + obstacle avoidance}
    \gbar{2}{7}{12}{cver}{VRX LiDAR/TF fix}
    \gbar{3}{13}{16}{cver}{Node rename + parameter sync}
    \gbar{4}{17}{21}{cver}{Dead-code audit + hardening}
    \gbar{5}{15}{21}{cpre}{Hardware-transition prep}
    \gbar{6}{24}{27}{cver}{Presentation + scope lock}
    \gbar{7}{31}{38}{cver}{RTF fix, CSP, first wet test}
    \gbar{8}{38}{47}{cver}{Pi 5 bring-up + doc audit}
    \gbar{9}{48}{56}{cpre}{Pi reflash + MAVROS install}
    \gbar{10}{59}{77}{cver}{FCU telemetry + QGC bridge}
    \gbar{11}{80}{88}{cver}{Onboard perception + thermal}
    \gbar{12}{89}{98}{cver}{Hailo-8L bring-up}
    \gbar{13}{101}{112}{cver}{Live view-only dashboard}
    \gbar{14}{115}{122}{cpre}{FCU-to-VRX bridge + prep}

    % ---- legend -------------------------------------------------------------
    \fill[cver] (0,15.6) rectangle (7,16.0);
    \node[anchor=west, black!85] at (9,15.8) {Verified};
    \fill[cpre] (40,15.6) rectangle (47,16.0);
    \node[anchor=west, black!85] at (49,15.8) {Prepared, not run};
  \end{tikzpicture}
  \caption{Overall project schedule, 03/04/2026 to 03/08/2026, colour-coded by evidence state.}
  \label{fig:gantt}
\end{figure}

\begin{longtable}{@{}L{0.20\textwidth}L{0.44\textwidth}L{0.14\textwidth}@{}}
  \caption{Overall Professional Thesis project schedule}\label{tab:project-schedule}\\
  \toprule
  \textbf{Period} & \textbf{Phase} & \textbf{State} \\
  \midrule
  \endfirsthead
  \caption[]{Overall Professional Thesis project schedule (continued)}\\
  \toprule
  \textbf{Period} & \textbf{Phase} & \textbf{State} \\
  \midrule
  \endhead
  03/04--09/04/2026 & Code review and obstacle-avoidance hardening & Verified \\
  10/04--15/04/2026 & VRX LiDAR/TF fix and supervisor review & Verified \\
  16/04--19/04/2026 & Node rename and parameter single-source sync & Verified \\
  20/04--24/04/2026 & Dead-code audit and reliability hardening & Verified \\
  18/04--24/04/2026 & Hardware-transition preparation (parallel) & Prepared \\
  27/04--30/04/2026 & Presentation preparation and scope lock & Verified \\
  04/05--11/05/2026 & Real-time-factor fix, CSP hardening, first wet test & Verified \\
  11/05--20/05/2026 & Raspberry Pi 5 bring-up and documentation audit & Verified \\
  21/05--29/05/2026 & Pi reflash, MAVROS install, endpoint gate & Prepared \\
  01/06--19/06/2026 & Flight-controller telemetry and mission bridge & Verified \\
  22/06--30/06/2026 & Onboard perception and thermal characterisation & Verified \\
  01/07--10/07/2026 & Hailo-8L bring-up and detector recovery & Verified \\
  13/07--24/07/2026 & Live detection overlay and view-only dashboard & Verified \\
  27/07--03/08/2026 & Autopilot-to-simulator bridge and query hardening & Prepared \\
  \bottomrule
\end{longtable}

An earlier project phase, from 25/11/2025 to 14/12/2025, established the initial architecture and
autonomous-navigation baseline. It is recorded as context and is not counted as internship delivery.

\ReportTodo{Add the post-03/08/2026 period once the internship closes.}

\clearpage
\pagenumbering{arabic}

% Length target from the instructions: 40 pages +/- 5, EXCLUDING appendices,
% including all text, figures, tables, and diagrams. The instructions do not
% say whether front matter counts; confirm with the internship supervisor.

% == INTRODUCTION =============================================================
\chapter*{Introduction}
\addcontentsline{toc}{chapter}{Introduction}

\section*{Subject and problem statement}

Monitoring water quality in ports, canals, and inland water bodies commonly depends on manual
sampling at a limited set of locations. Published unmanned-surface-vehicle systems show that
autonomous navigation, obstacle avoidance, water-quality sensing, sample collection, and remote
operator displays can be integrated on one mobile platform~\cite{chang2021waterquality}. The
engineering problem addressed by this Professional Thesis is the one that stands between that idea
and a working system: building an autonomous navigation and perception stack that behaves
predictably in simulation, and then transferring it onto real marine hardware without losing the
guarantees established in simulation.

\ReportTodo{Add why this problem matters specifically to the host institution --- the operational
need, the users, and what existed before the project.}

\section*{Professional Thesis objectives}

The work was organised around four objectives:

\begin{enumerate}
  \item Make the simulated navigation stack reliable and measurable, so that behaviour can be
        compared before and after a change rather than judged by eye.
  \item Establish a verified telemetry link between a companion computer and the vessel's flight
        controller, as the foundation of any later command capability.
  \item Bring onboard perception up on an edge AI accelerator, and prove the detection path end to
        end from camera to operator display.
  \item Provide an operator-facing dashboard that presents live detection and vehicle telemetry
        together, under an explicit read-only safety contract.
\end{enumerate}

A fifth, research-oriented objective --- water-quality sensor streaming, a cellular-automaton
evolution model, and machine-learning trend detection --- was defined but is reported here as not
started, for the reasons set out in the Conclusion.

\section*{Monitoring and reporting plan}

Work was tracked in a dated engineering diary covering 03/04/2026 to 03/08/2026, one file per
working day, recording what was attempted, what was measured, and what remained open. Each
significant change was verified before being recorded as complete, and every claim in this report
is tagged by evidence state:

\begin{itemize}
  \item \textbf{Verified} --- executed and observed, with a measured result.
  \item \textbf{Prepared} --- implemented and statically checked, but never executed.
  \item \textbf{Unrun} --- planned, scaffolded, or blocked by an external dependency.
\end{itemize}

This distinction is used consistently and deliberately. Several parts of the system are complete in
source but have never been run against real hardware, and presenting them as results would
misrepresent the state of the project.

\section*{Report structure}

Section~\ref{chap:context} places the assignment in its institutional and sector context.
Section~\ref{chap:project} presents the assignment itself, broken into five sub-projects, each
following a diagnose--solve--operationalize--evaluate structure, and closes with the consolidated
results. Section~\ref{chap:skills-career} assesses the skills developed and their relationship to a
professional plan.

% == SECTION 1 - CONTEXT ======================================================
% Mandatory. MAXIMUM 5 PAGES.
% Must not be filler, plagiarism of company documents, or reproduced website text.
\chapter{Context and Institutional Expectations}
\label{chap:context}

\section{Host institution and wider sector}
Within the wider sector, mobile robotic platforms offer spatially distributed measurements while
also carrying the navigation, communications, and operator-interface functions needed to associate
each observation with a location~\cite{chang2021waterquality}. This project addresses that enabling
navigation, perception, telemetry, and operator-monitoring layer. Water-quality sensor streaming,
spatio-temporal modelling, and trend detection remain future work rather than results of the present
thesis.

\ReportTodo{Present the institution without copying its website: its position, its activity in this
sector, and how the autonomous-vehicle work fits its wider objectives. This cannot be drawn from
the project repository and must be written from direct knowledge.}

\PlaceholderFigure{org-context}{Position of the project's team or laboratory within the host
institution, and its relationship to the wider autonomous-systems sector.}{Organisational context
of the Professional Thesis assignment.}

\section{Organization and working environment}
\ReportTodo{Describe only the locations, organization, processes, and working methods needed to
understand the assignment --- for example the balance between individual work and supervision, and
the equipment and workshop access the project depended on.}

\section{Strategic and technical challenges}

The technical environment set three constraints that shaped every design decision, and these are
evidenced by the project record.

First, the vessel is a differential-thrust catamaran: it has no rudder, and all steering comes from
commanding unequal power to left and right thrusters. Every navigation decision therefore resolves
into a thrust pair, which makes the controller the single point through which all planning must
pass.

Second, the onboard computing budget is small. A Raspberry Pi 5 companion computer carries the
high-level autonomy alongside a separate autopilot board, with a PCIe neural accelerator for
inference. Thermal and power headroom are real limits rather than theoretical ones; undervoltage
warnings and accelerator temperature were both monitored during live runs.

Third, the operator link is shared and constrained. Manual control, telemetry, and video all pass
through an integrated radio system, and the workstation, vessel, and radio do not always sit on the
same network. Several field sessions were spent establishing which network could carry which
traffic.

\ReportTodo{Add the institution's own strategic framing --- the digital or ecological transition,
research programme, or funding context this project sits inside.}

\section{Project importance and institutional expectations}
\ReportTodo{State what the institution expected from this Professional Thesis, what would count as
success at the outset, and how that expectation was formalised.}

\section{Personal assessment of the institution's SDRS policy}
% Explicitly required by the instructions - do not drop this section.
\ReportTodo{Give a substantiated personal opinion on the institution's sustainable-development and
social-responsibility policy. Confirm the official terminology used locally, and support the
assessment with what was observed during the internship rather than with published material.
A defensible angle available here: the project's own environmental purpose --- water-quality
monitoring --- against how sustainability is handled in day-to-day engineering practice.}

% == SECTION 2 - PROJECT ======================================================
% Mandatory. No page limit stated - this is the core of the report.
\chapter{Objectives and Professional Thesis Project}
\label{chap:project}

\section{Mission, objectives, and expected value}

The assignment was to advance an autonomous surface vehicle from a simulation-only prototype toward
a field-capable system. The deliverables were a reliable simulated navigation stack, a verified
telemetry path to the real vessel, onboard perception running on an edge accelerator, and an
operator dashboard combining both. The value to the institution is a platform on which
water-quality survey missions can later be built, together with a documented evidence trail
showing exactly which capabilities are proven and which are not.

Figure~\ref{fig:system-arch} summarizes the active ROS 2 data flow and its external interfaces.

\PlaceholderFigure{system-arch}{ROS 2 node pipeline: sensors to perception to planning to control
to differential thrust, with topic names on each arrow, the latched emergency-stop channel, and
the two configuration services.}{System architecture of the autonomous navigation stack.}

\section{Resources and constraints}

The material resources were a simulation workstation with a discrete GPU, the vessel with its
autopilot board and integrated radio, a companion computer with a neural accelerator and a depth
camera, and access to open water for field tests. Human resources were the author, working with
academic supervision and periodic review. The dominant constraints were not financial but
practical: single-owner hardware devices that cannot be shared between processes, a thermal ceiling
on the companion computer, a shared radio link, and field sessions that had to be scheduled and
could not be repeated cheaply.

Two constraints were treated as absolute safety boundaries throughout. No command was ever sent to
a real motor, and no arming of the real vessel was attempted outside a restrained bench
configuration. These were enforced in software, not only by convention, as described in
Section~\ref{sec:sub-dashboard}.

Figure~\ref{fig:deployment} separates the system components by host and communication boundary.

\PlaceholderFigure{deployment}{Deployment topology: autopilot board, companion computer, integrated
radio, and workstation, showing which wireless network carries which traffic and marking the
receive-only serial link.}{Field deployment topology across the three wireless networks.}

% ---------------------------------------------------------------------------
\section{Sub-project 1: bounding detour-waypoint growth in obstacle avoidance}
\label{sec:sub-avoidance}

\subsection{Diagnose}

During simulated missions the vehicle inserted more than 300 detour waypoints on encountering an
obstacle and still failed to bypass it; a separate planning run grew the waypoint list from 203 to
over 303 entries. Analysis identified seven contributing defects across two nodes. On the planner
side, the obstacle callback reset its ``detour inserted'' flag whenever obstacles momentarily
cleared between scans, permitting immediate re-insertion; there was no cumulative cap and no
cooldown; the insertion routine always chose the left side without checking clearance; and no path
validated that the detour point was itself clear. On the controller side, the maximum avoidance
turn was limited to 20 degrees, the escape manoeuvre always turned left, and a replanning method
was called but never defined, producing a silent failure during high-urgency encounters.

\subsection{Solve}

The chosen approach bounded growth rather than redesigning the planner, on the reasoning that an
unbounded feedback loop is a correctness fault and should be fixed as such before any behavioural
improvement is attempted. The premature flag reset was removed; a counter capping detours at three
per waypoint and a five-second cooldown were added; a clearance validator was introduced that tries
the chosen side, then the opposite side, then skips the waypoint; and the side is now selected from
measured left and right clearance instead of a fixed preference. The missing replanning method was
defined, the maximum avoidance turn was raised from 20 to 45 degrees scaled by urgency, and the
escape manoeuvre now turns toward the side with greater measured clearance. A burst-skip mechanism
was added so that a long linear obstacle such as a pier produces one skip past the blocked stretch
rather than three detours per waypoint.

\subsection{Operationalize}

All three detour insertion paths were made consistently capped, including one that previously had
no cap at all, and code defaults were synchronised with the launch configuration. The new avoidance
limit was added to the automated health-check script so that a regression would be caught by the
routine monitoring run rather than by observation. Dashboard tuning presets and the corresponding
interface defaults were re-synchronised, so that applying a preset could no longer silently
override the new values. Dead code found in the same pass was removed.

\subsection{Evaluate}

The post-change health check returned a full pass with no failures or warnings. An open-water
mission completed 18 waypoints in 4.9 minutes, with the obstacle encountered at approximately
5.6~m, urgency peaking at 97~\%, two detours inserted, no runtime error, and the total waypoint
count remaining at 18. A pier scenario completed in 9.5 minutes with the count bounded at 17 rather
than exceeding 300. The record is explicit that the problem was bounded rather than solved: the
vehicle still does not bypass long piers smoothly, and pier and bank stuck-behaviour remained the
top open navigation issue at the end of the period.

Figure~\ref{fig:detour-growth} contrasts the original feedback growth with the bounded result.

\PlaceholderFigure{detour-growth}{Detour-waypoint count before and after the fix: unbounded growth
past 300 insertions, against the bounded post-fix behaviour of 18 and 17 waypoints.}{Detour
waypoint growth before and after bounding.}

% ---------------------------------------------------------------------------
\section{Sub-project 2: root-causing the simulation real-time-factor throttle}
\label{sec:sub-rtf}

\subsection{Diagnose}

The simulation ran at roughly one fifth of real time, which slowed every test and made
hardware-versus-simulation comparison unreliable. Baseline measurement recorded the LiDAR point
cloud at 2.48~Hz against a 10~Hz nominal rate, and the simulation clock at 80.9~Hz against a 250~Hz
nominal, giving a real-time factor of 0.32. The LiDAR shortfall was proportionally worse than the
overall ratio, which pointed at the GPU-bound raycasting sensor rather than at a general CPU
shortage. The prior working hypothesis, a graphics-library fallback, was disproven: the discrete
GPU driver loaded cleanly, but the distribution's default hybrid-graphics mode meant the launcher
never requested offload, so the simulator rendered on the integrated GPU.

\subsection{Solve}

Five ways of applying the offload were compared: a per-invocation environment prefix, a system-wide
switch, unconditional baking of the offload variables into the launcher, automatic detection of a
discrete GPU, and an opt-in launcher flag. Unconditional baking was rejected because it would force
the variables onto hosts with no discrete GPU, notably the companion computer; automatic detection
was rejected as less predictable for a reader of the script. The opt-in flag was chosen on the
reasoning that on the companion computer an accidental invocation then fails visibly at graphics
context creation rather than degrading silently --- a failure mode preferred deliberately.

\subsection{Operationalize}

The flag shipped together with a launcher signal-handling fix, so that the script no longer survived
an interrupt. The troubleshooting documentation was rewritten as a three-step diagnostic flow
carrying the measured before-and-after figures and a hybrid-graphics-only caveat, retaining the
manual environment-prefix form as a fallback for one-off launches. The flag was then surfaced in
five user-facing documents, because a reader following the main guide would otherwise land in
integrated-GPU mode by default.

\subsection{Evaluate}

Two measured runs were compared against the baseline. Standalone simulation reached LiDAR 8.95~Hz,
clock 195.9~Hz and a real-time factor of 0.78; the full launcher reached LiDAR 6.8~Hz, clock
219.7~Hz, factor 0.88, and a 44-second launch time. This is a 2.7-fold improvement in
full-launcher LiDAR rate and 2.4 to 2.7-fold in real-time factor. Workstation thermals during the
accelerated run stayed well below the throttle ceiling. A secondary contribution from CPU frequency
governance was estimated at a further 10 to 15~\% but was never measured, and is reported as an
estimate only.

Figure~\ref{fig:rtf-improvement} compares the measured real-time factor and LiDAR rate across the
recorded graphics configurations.

\PlaceholderFigure{rtf-improvement}{Simulation performance across three configurations: integrated
GPU baseline, discrete standalone, and discrete full launcher, plotting LiDAR rate against the
10~Hz nominal and real-time factor against 1.0.}{Simulation performance before and after the
graphics-offload fix.}

% ---------------------------------------------------------------------------
\section{Sub-project 3: establishing a telemetry link to the flight controller}
\label{sec:sub-telemetry}

\subsection{Diagnose}

At the start of June no usable communication endpoint existed on the companion computer. A serial
sweep found only a debug-class device, with none of the expected serial or USB device nodes and no
network listener on any of the standard ports. Three independent tools --- the middleware bridge, a
protocol router, and a direct protocol script --- were then used as cross-checks at three baud
rates. All three opened the port and all three received nothing; a raw dump of the device showed
zero bytes. Since a baud mismatch normally produces garbled bytes rather than silence, the fault was
reasoned to be physical or configuration-side rather than a software fault in the bridge. A
candidate service was removed and the machine rebooted with no change, retiring it as an
explanation and leaving the platform's serial-port mapping as the leading suspect.

\subsection{Solve}

Candidate endpoint classes were enumerated and scored on available evidence: USB, unusable because
no device ever enumerated; network, unusable because listeners opened but no source was confirmed;
and a hardware serial connection to the general-purpose header requiring a platform-specific
configuration change. The third branch was confirmed by a supervisor-supplied photograph showing the
protocol router attached to the expected device at 57\,600 baud and reporting a detected vehicle. A
second architectural choice followed: binding the middleware bridge directly to the serial port, or
letting the router own the serial link and redistribute over the local network. The redistribution
topology was selected because a single serial owner can feed the bridge plus optional inspection
tools without device contention.

\subsection{Operationalize}

The proven topology is the protocol router as sole owner of the serial device, redistributing to two
local network ports: one consumed by the middleware bridge, the other reserved strictly for direct
inspection scripts and never for the bridge. A shared domain identifier was required in every
terminal on both machines and was later pinned into the shell profile on each. Monitoring used a
kernel-log undervoltage check before and after each run, a connection-state polling loop as the
session gate, and a best-effort delivery profile on sensor telemetry. The gate rule was written
explicitly, and is the methodological point of this sub-project: advertised topics never count as
success, only a confirmed connection state does.

\subsection{Evaluate}

The pass condition was reached on 04/06/2026 with the link reporting connected, disarmed, and in
hold mode, and first samples arriving on the inertial, position, battery, and radio-input topics. A
clean graph capture the following day contained 144 typed topics, of which 136 belonged to the
bridge. A re-check on 19/06/2026 after a package synchronisation passed with no regression. Three
limitations are preserved in the record: satellite positioning remained at no-fix throughout;
targeted request-and-response paths timed out, so the command and write path is explicitly
unvalidated; and a combined camera-plus-telemetry run was power-limited under repeated undervoltage
warnings.

Figure~\ref{fig:telemetry-chain} traces the verified read path from the flight controller to the
observed middleware topics.

\PlaceholderFigure{telemetry-chain}{Telemetry chain from flight controller to middleware topics:
controller telemetry port, companion-computer serial device, protocol router as sole owner, local
network redistribution, and the resulting topic namespace, annotated with the connection gate and
marking the write path unvalidated.}{Verified telemetry chain from the flight controller.}

% ---------------------------------------------------------------------------
\section{Sub-project 4: live detection overlay and telemetry in one dashboard}
\label{sec:sub-dashboard}

\subsection{Diagnose}

Three gaps blocked a combined live view. The inference runner was an external process that owned the
camera device directly and wrote annotated frames to a file only, with no topic and no network
output, whereas the dashboard camera panel is a plain image element fed by a streaming server. The
depth camera is single-owner, so the camera driver and the inference runner cannot co-run, meaning
annotated frames had to originate inside the runner. The dashboard also subscribed only to
simulation topics and had no inertial panel, so the originally written acceptance criterion had to
be replaced. Operationally, early runs were launched by hand across many terminals and left no
same-run evidence: observations were real but produced no completion or teardown markers.

\subsection{Solve}

Five transport options were compared for the video path. A thin publishing node inside the runner,
emitting the annotated frame on a standard image topic, was chosen because it matches the
dashboard's existing topic-discovery filter with no dashboard edit, reuses the already-proven
cross-host transport, and exposes no new unauthenticated port on the vessel. Serving a stream
directly from the companion computer was rejected because it requires editing a hard-coded host and
opens an unauthenticated port; a streaming-protocol option was rejected because it needs a transcode
hop the image element cannot render. For telemetry, an external read-only relay onto the dashboard's
existing topics was preferred, with the thrust and command direction deliberately left unmapped. For
lifecycle, the existing preflight script was extended with a supervisor mode rather than writing a
new supervisor, and a hold-never-abandon failure model was chosen over exiting and leaving orphaned
processes.

\subsection{Operationalize}

The stack was reduced to two operational commands: a foreground workstation supervisor that runs a
preflight, starts each service in its own process group with external logs, prints one
checksum-pinned copy-paste command for the vessel, blocks on a six-topic arrival gate, runs rate
probes automatically, supervises, and tears down in reverse order; and a vessel-side helper owning
the serial device, the router, the bridge, the inference runner, and the camera. Safety and evidence
were built into the contract rather than left to discipline: a view-only flag blocks every dashboard
publish and service call at the send boundary; a safety monitor subscribes to all five protected
command topics and aborts the run on any traffic or on arm or disconnect; a one-second thermal
watchdog aborts at a fixed ceiling; binary and helper checksums are verified at launch; and a
durable per-run log records phase, first stop trigger, and final status.

\subsection{Evaluate}

Two bounded runs on 17/07/2026 proved the combined path: all six expected topics arrived after 248
and 260 seconds, the detection overlay sustained 7.40 and 7.50~Hz with five telemetry topics at
approximately 1.00~Hz, the link reported connected and disarmed, both command sentinels closed with
zero messages on all five protected topics, and thermal peaks stayed below the abort ceiling.
Operator visual acceptance passed for both runs and is reported separately from the rate evidence. A
later run reached the arrival gate in 111 seconds, and a dual-output run showed the annotated view
simultaneously on the vessel display and the workstation dashboard with at least 2\,800 annotated
publications. Open items are recorded explicitly: a strict shutdown ordering was not achieved and
its cause remains open; endurance acceptance and post-teardown thermal behaviour remain open; and no
dashboard-to-controller write path was ever exercised.

The browser-facing path uses the rosbridge JSON interface~\cite{rosbridge}.
Figure~\ref{fig:live-dataflow} distinguishes the live view-only data path from the protected command
topics.

\PlaceholderFigure{live-dataflow}{Data paths of the live view-only test: camera to accelerator to
image topic to streaming server to browser, alongside controller to router to bridge to dashboard
telemetry panels, with the view-only guard and command sentinel drawn as the safety
boundary.}{Data paths and safety boundary of the live view-only test.}

% ---------------------------------------------------------------------------
\section{Sub-project 5: confirming and routing around the command path}
\label{sec:sub-command}

\subsection{Diagnose}

A read-only audit established that no software path currently reaches a physical motor. The thrust
topics are published by the heading controller and consumed only by the simulator plugin; on real
hardware nothing subscribes to them. The real-hardware translator in the launch configuration is an
unbuilt stub, and requesting real hardware aborts the launch. The middleware bridge runs a bounded
telemetry allowlist with no command, setpoint, or actuator plugins loaded, and the dashboard
view-only guard blocks every publish at the send boundary. A separate link probe then found that the
companion computer can receive from the controller but cannot transmit to it: parameter reads and
command requests returned no acknowledgement, a motor test left both servo channels at neutral,
while arming from the handheld ground station succeeded.

\subsection{Solve}

The real vehicle interface was identified from stored ground-station telemetry rather than assumed:
the thruster outputs sit on the servo rail with a known channel and function mapping, a PWM range of
800 to 2200 with neutral at 800, and arming required together with a hardware safety switch. The
conclusion recorded was that the only real command path requires arming, which is a stated non-goal
and a safety-monitor abort trigger. Two hardware suspects were named without being asserted as
causes: an unconnected transmit line, or flow control on a three-wire link. A route around the
blocker was then identified: a network-based bridge means that against a simulated autopilot on the
loopback interface nothing crosses the blocked serial link, making the autopilot-to-simulator
command path developable while the wiring fix is pending.

\subsection{Operationalize}

A reference implementation was reviewed and four adaptation issues found, the most serious being a
neutral-point assumption that, applied unchanged to this vessel, would drive the simulator at
roughly 80~\% forward while the autopilot commanded stop. A runnable equivalent was then written as
a single standalone file requiring no new package, with deliberate safety-driven differences: it
identifies as an onboard controller and never impersonates the vehicle, decodes this vessel's real
channel mapping and PWM range, publishes the project's own thrust topics, keeps sensor injection
toward the autopilot disabled by default because transmitting into a real controller is a write, and
uses bounded receive with clean shutdown. Its boundary is documented: it publishes on topics the
safety monitor treats as protected, so it must run in simulation on an isolated domain and never
beside the live vessel stack.

\subsection{Evaluate}

No live motor command was sent, and no dashboard write path, view-only guard, or safety monitor was
changed. The bridge is source-verified and harness-tested only. A local harness with synthetic
protocol input verified the safety lifecycle: interrupt, termination, and repeated signals each
exited cleanly and published a trailing zero on both thrust topics after held non-zero samples, and
configuration validation rejected non-finite or non-positive thrust limits, duplicate channels, and
an out-of-range neutral value. The integrated simulated-autopilot run is recorded as \textbf{not
run}, the PWM starting values are provisional rather than confirmed, and the installation and launch
commands are labelled untested. The track is parked behind two explicit gates: a bidirectional link,
with success defined as reading a parameter back from the controller, and a separate
propellers-removed bench-safety decision.

% ---------------------------------------------------------------------------
\section{Autonomy, responsibility, and decision-making}
Three technical decisions demonstrate initiative in the recorded work. The obstacle-avoidance
feedback loop was bounded before further behavioural tuning (Section~\ref{sec:sub-avoidance});
three of five data-transport options were rejected because they weakened safety or maintainability
(Section~\ref{sec:sub-dashboard}); and telemetry acceptance was defined by observed runtime
messages rather than advertised topic names (Section~\ref{sec:sub-telemetry}). Hardware-facing
work remained approval-gated.

\ReportTodo{State which of these decisions were made independently and which were validated with a
supervisor, then describe the corresponding level of responsibility.}

\section{Consolidated results and company benefits}

Table~\ref{tab:indicators} consolidates the measured indicators. Every value is quoted as recorded
at the time of measurement, and each is tagged by evidence state.

\begin{table}[tbp]
  \caption{Consolidated performance indicators.}
  \label{tab:indicators}
  \centering
  \small
  \begin{tabularx}{\textwidth}{@{}L{0.30\textwidth}YL{0.11\textwidth}@{}}
    \toprule
    \textbf{Indicator} & \textbf{Measured value} & \textbf{State} \\
    \midrule
    Detour-waypoint growth & Before: over 300 inserted. After: 18 waypoints in 4.9~min (open
      water), 17 (pier scenario) & Verified \\
    Automated health check & 49 runtime checks; 49 pass idle, 41 pass and 8 tuned active, 0 fail
      & Verified \\
    Simulation real-time factor & 0.32 to 0.88; LiDAR 2.48 to 6.8~Hz; clock 80.9 to 219.7~Hz
      & Verified \\
    Perception publish rate & 20.000~Hz mean over 2\,353 samples in 120~s, no sustained dropouts
      & Verified \\
    Simulated mission completion & 16 of 16 waypoints, 405~s simulated time & Verified \\
    Web security hardening & 0 policy violations across 11 interactive paths, 2 browsers; injected
      payloads rendered as literal text & Verified \\
    Cross-machine discovery & 162 messages delivered; latency 7.89 to 11.8~ms & Verified \\
    Telemetry link gate & Connected, disarmed, hold mode; 136 bridge topics on the shared domain
      & Verified \\
    Mission bridge to ground station & 7 mission items served and rendered without file import;
      15 of 15 unit tests pass & Verified \\
    Accelerator runtime & 293 frames at 58.22~FPS; 287 frames at 57.29~FPS on an independent boot
      & Verified \\
    Host-side decode contract & Same-engine isolation: box error 0.0~px, class error
      $1.178\times10^{-7}$ & Verified \\
    Live combined overlay run & Overlay 7.40 and 7.50~Hz; telemetry near 1.00~Hz; 0 messages on all
      five protected command topics & Verified \\
    \bottomrule
  \end{tabularx}
\end{table}

One negative result belongs beside the table, because omitting it would misrepresent the state of
perception. The custom detector fires on none of the 50 available saved frames at the working
confidence threshold, including on its own training images, with a confidence ceiling near 0.003.
The underlying cause is a dataset problem rather than a pipeline problem: the baseline dataset holds
nine labelled instances of a single class, with four declared classes carrying no examples at all.
The inference pipeline itself is proven, since the host-side decode contract matches to within
$1.178\times10^{-7}$ and a stock model produces correct overlays at full frame rate.

Figure~\ref{fig:detector-collapse} summarizes the detector collapse across thresholds and data
splits.

\PlaceholderFigure{detector-collapse}{Detection count and maximum confidence at three thresholds
for the validation split against the held-out split, annotated with the training and held-out
scores and the resolution sweep.}{Held-out firing collapse in the detector training smoke test.}

Figure~\ref{fig:dashboard} identifies the operator-facing elements exercised during the bounded
view-only runs.

\PlaceholderFigure{dashboard}{The operator dashboard in use: mission control, map, live camera
panel, and health-check stream, with the emergency-stop badge and configuration
panels.}{Operator dashboard during a live view-only run.}

The technical benefit is a reproducible near-real-time simulation baseline, a verified read path
from the flight controller, and bounded operation of the combined overlay and dashboard under an
enforced view-only contract. Open limitations are attached to explicit acceptance gates, so later
work can start from observed evidence rather than assumed capability. Command-path acceptance,
detector recovery, and full lifecycle endurance remain open.

\ReportTodo{Relate these capabilities to the institution's own priorities and quantify their
operational value using information supplied by the host team.}

% == SECTION 3 - SKILLS AND CAREER ============================================
% Mandatory. MAXIMUM 5 PAGES.
\chapter{Skills Development and Professional Future}
\label{chap:skills-career}

\section{Skills assessment}
\ReportTodo{Assess the skills developed using the IMT Nord Europe engineering skills toolkit, naming
the toolkit's own competence categories. Support each with evidence from
Section~\ref{chap:project}. The record supports at least: systematic fault diagnosis under
uncertainty; designing for verifiability rather than for demonstration; and safety-boundary design
in a system able to move physical hardware.}

\section{Functional and managerial observations}
\ReportTodo{Analyze relationships between teams, distribution of roles, management, working
conditions, safety practice, sustainable development, social responsibility, communication, and
institutional culture, based on what was observed rather than on published material.}

\section{Development needs and further training}
Two technical development priorities follow directly from the recorded blockers. First, the
held-out detector collapse shows a need for stronger dataset design, annotation quality control,
and validation methodology. Second, the receive-only serial link shows a need for deeper
embedded-hardware bring-up skills, including signal-level diagnosis and flight-controller serial
configuration.

\ReportTodo{State the personal training plan for these priorities and whether further formal study
is intended.}

\section{Professional objectives}
\ReportTodo{Connect the experience to the sector, target roles, and the next career step, and
explain how this project supports that transition.}

% == CONCLUSION ===============================================================
\chapter*{Conclusion}
\addcontentsline{toc}{chapter}{Conclusion}

The problem set at the outset was to move an autonomous surface vehicle from a simulation prototype
toward a field-capable system without losing, in the transfer, the guarantees established in
simulation. Against that problem the outcome is partial and asymmetric, and it is worth stating
precisely which way.

The simulation stack and the read side of the real system are proven. Obstacle-avoidance growth is
bounded and regression-guarded; simulation runs near real time after the graphics-offload fix; a
full mission completes; perception publishes at a stable rate; the web interface has no policy
violations; a telemetry link to the real flight controller reaches a confirmed connected state and
sustains 136 topics; the neural accelerator runs at full frame rate with a decode contract matching
to within $1.178\times10^{-7}$; and two bounded live runs displayed detection overlay and vehicle
telemetry together in a browser under an enforced read-only contract.

The write side is not proven, and no amount of prepared work changes that. The link from companion
computer to flight controller is confirmed receive-only, so no actuator, arming, mode, or thrust
command has ever been attempted; the two hardware suspects remain suspects rather than diagnosed
causes. The autopilot-to-simulator bridge is source-verified and harness-tested for safe teardown,
but its integrated run has not been performed. The hardware translator package does not exist. The
custom detector does not fire, and the cause is an insufficient dataset rather than a pipeline
fault. The entire research layer --- sensor streaming, the evolution model, and trend detection ---
is not started. There has been no autonomous on-water run; the single field session was a
manual-control float test.

The principal professional lesson concerns evidence rather than technique. Several times during the
period a component was complete in source and felt finished, and each time the discipline that
mattered was refusing to record it as a result until it had been run and observed --- and, more
than once, discovering on running it that the confident version had been wrong. The habit of
separating verified from prepared from unrun, and of writing acceptance criteria that cannot pass
vacuously, was the most transferable thing built during this Professional Thesis.

\ReportTodo{Close with a short forward statement once the internship ends: what was completed after
03/08/2026, and what is handed over.}

% == TECHNICAL GLOSSARY =======================================================
\chapter*{Technical Glossary}
\addcontentsline{toc}{chapter}{Technical Glossary}

Table~\ref{tab:technical-glossary} defines the technical terms and acronyms used in the report.

\begin{longtable}{@{}L{0.25\textwidth}L{0.67\textwidth}@{}}
  \caption{Technical glossary}\label{tab:technical-glossary}\\
  \toprule
  \textbf{Term or acronym} & \textbf{Definition in this report} \\
  \midrule
  \endfirsthead
  \caption[]{Technical glossary (continued)}\\
  \toprule
  \textbf{Term or acronym} & \textbf{Definition in this report} \\
  \midrule
  \endhead
  USV & Unmanned surface vehicle: an uncrewed boat operating on the water surface. \\
  ROS 2 & The robotics middleware used throughout~\cite{macenski2022ros2}; a \emph{node} is one
    process, a \emph{topic} a named typed message channel. \\
  DDS / domain ID & The peer-to-peer discovery layer beneath ROS 2, and the numeric namespace that
    isolates one system's traffic on a shared network. \\
  QoS & Per-topic delivery policy; incompatible publisher and subscriber settings exchange nothing,
    silently. \\
  VRX / WAM-V & The maritime simulation suite~\cite{bingham2019vrx} and the catamaran-style vessel
    it simulates, running on Gazebo~\cite{gazebo}. \\
  Real-time factor & Ratio of simulated to wall-clock time; every simulator-measured rate must be
    qualified by it. \\
  LiDAR / point cloud & Laser range sensor producing roughly 30\,000 3D points per scan. \\
  Differential thrust & Steering by unequal left and right thruster power rather than by rudder;
    the vessel's only steering mechanism. \\
  Emergency stop & A latched Boolean on a dedicated reliable-delivery topic, subscribed
    independently by planner and controller. \\
  MAVLink & The binary protocol~\cite{mavlink} spoken between autopilots, ground stations, and
    companion computers; bridged into ROS 2 by MAVROS~\cite{mavros}. \\
  Flight controller & The autopilot board driving the thrusters, running ArduPilot's rover
    firmware~\cite{ardupilot}; distinct from the companion computer. \\
  Companion computer & The onboard general-purpose Linux computer running high-level autonomy
    alongside the autopilot. \\
  Edge AI accelerator & The onboard neural inference device~\cite{hailo}; its binary format exposes
    raw outputs with no embedded post-processing. \\
  Object detection & Locating and classifying objects in an image, here with a single-stage
    detector family~\cite{ultralytics}. \\
  Software in the loop & Running real autopilot firmware as a simulated vehicle on an ordinary
    computer, over network sockets rather than a serial link. \\
  Digital twin & The simulator kept aligned with the real vessel by streamed telemetry. \\
  Cellular automaton & A grid model in which each cell's next state depends on its neighbours;
    planned, not implemented. \\
  \bottomrule
\end{longtable}

% == INFORMATION SOURCES ======================================================
% Every non-original figure or table must cite its source from here.
\bibliographystyle{plain}
\bibliography{references}

\ReportTodo{Before final submission, verify and add primary references for the path-planning and
detector methods where they are discussed, and ensure every non-original illustration cites its
source.}

% == APPENDICES ===============================================================
% Appendices are EXCLUDED from the 40-page count and must be numbered.
\appendix

\chapter{Terminology Appendix}
\label{app:terminology}

Because this report is written in English, this appendix provides French equivalents for specialized
terms needed to interpret the technical discussion.

Table~\ref{tab:terminology} records the corresponding English and French terminology.

\begin{longtable}{@{}L{0.28\textwidth}L{0.28\textwidth}L{0.36\textwidth}@{}}
  \caption{Terminology appendix: specialized terms and translations}\label{tab:terminology}\\
  \toprule
  \textbf{Report term} & \textbf{Translation} & \textbf{Usage note} \\
  \midrule
  \endfirsthead
  \caption[]{Terminology appendix (continued)}\\
  \toprule
  \textbf{Report term} & \textbf{Translation} & \textbf{Usage note} \\
  \midrule
  \endhead
  Unmanned surface vehicle (USV) & véhicule de surface autonome &
    Autonomous boat operating at the water surface. \\
  Digital twin & jumeau numérique &
    Simulation aligned with the physical vessel through telemetry. \\
  Telemetry & télémétrie &
    Status data received from the flight controller. \\
  Object detection & détection d'objets &
    Image-based localization and classification of objects. \\
  \bottomrule
\end{longtable}

\end{document}
